\documentclass[10pt,a4paper]{amsart}
\usepackage[margin=19mm]{geometry}
\usepackage{amsmath,amssymb,mathtools}
\usepackage[scr=rsfs,cal=euler]{mathalfa}
\usepackage{xfrac}
\usepackage[unicode,hidelinks,bookmarks=false]{hyperref}
\newcommand{\ie}{i.e.\ }
\newcommand{\cf}{cf.\ }
\newcommand{\resp}{respectively\ }
\newcommand{\Subsection}{\subsection}
\providecommand{\nobreakdash}{\nobreak\hbox{-}}
\setlength{\emergencystretch}{2em}
\multlinegap0pt
\usepackage{bm}
\usepackage{bbm}
\usepackage{dsfont}
\usepackage{xspace}
\usepackage{cancel}
\usepackage{mathtools}
\usepackage{amsthm}
\makeatletter
\@ifundefined{lemma}{\newtheorem{lemma}{Lemma}}{}
\makeatother
\title[Minimizing Monochromatic Subgraphs of $K\_{n,n}$]{Minimizing Monochromatic Subgraphs of $K\_{n,n}$}
\author{Charles Gong}
\date{Source: arXiv:2410.19076v2 (CC BY 4.0)}
\begin{document}
\maketitle

\noindent\emph{Source and licence.} Minimizing Monochromatic Subgraphs of $K_{n,n}$. Version arXiv:2410.19076v2 (CC BY 4.0), distributed under the Creative Commons Attribution 4.0 International licence, as recorded in the arXiv licence metadata for this exact version and shown on the abstract page https://arxiv.org/abs/2410.19076v2. Full rights notice: licenses/arxiv-2410.19076-CC-BY-4.0.txt.\par\smallskip

\noindent\textit{Editorial scope.} Counted unit: the Lemma whose source label is \texttt{None} in the article (source0.tex lines 985--995, source label \texttt{None}), delivered together with the article's own complete original proof of it (source0.tex lines 997--1095, one \texttt{proof} environment of 99 lines). No mathematics is written, repaired or re-derived: statement and proof are the article's own text line for line. Required context retained uncounted: none. Every definition, display and label of those lines is retained unchanged. Omitted from this package and not redistributed: 1--984, 996--996, 1096--2207. Nothing inside a retained excerpt is omitted or altered: every retained body is delivered exactly as the article's lines are written, so no omission receipt of any kind accompanies this package. Citations: the retained excerpts carry no citation command at all, so no bibliography entry of the article is transcribed and no reference is redistributed. Reference adaptation: none was needed and none was made; every reference inside the delivered unit resolves inside it, so no unresolved reference or citation is produced. Printed slips kept literal: no printed word, symbol, hypothesis or number of the counted statement or of its proof was altered. Layout and class: the article's own class is \texttt{article} with packages including graphicx, amsmath, amsfonts, amssymb, babel, amsthm, tikz, thm-restate, geometry; the wrapper is the lane's pinned \texttt{amsart} editorial wrapper at 10pt on A4, which restates the theorem-like environments the retained text needs and adds the article's own macro definitions that the retained text needs (none), with extra wrapper packages (bm, bbm, dsfont, xspace, cancel, mathtools). The delivered numbering is the unit's own. Assets: no retained span contains a figure environment, a graphics-inclusion command, a PDF-page inclusion, a table or a TikZ picture; no image file is copied into this package, the package declares no asset dependency and no image file is redistributed with it. Licence: version arXiv:2410.19076v2 of this article is distributed under the Creative Commons Attribution 4.0 International licence, as recorded in the arXiv licence metadata for this exact version and shown on the abstract page https://arxiv.org/abs/2410.19076v2. A full rights notice is delivered at licenses/arxiv-2410.19076-CC-BY-4.0.txt. Characters: every retained excerpt is pure ASCII, so the delivered engine, XeLaTeX with its default Latin Modern text font, reports no missing character.
\begin{lemma}
Let $c_i, d_i \in [0,1]$ where $i \in [s]$, $k \in (0,1]$, and $[s] = I \sqcup O_1 \sqcup O_2$ where $|I| \geq 1$ and $|O_1| = |O_2| \geq 1$. If $c_i + d_i \leq k < \frac{2}{|I|+|O_1|}$ for all $i \in [s]$ and 

$$\sum_{i \in I \sqcup O_1} c_i \geq 1$$ 
$$\sum_{i \in I \sqcup O_2} d_i \geq 1,$$

\noindent
then if $e = \frac{2-|I| \cdot k}{2|O_1|}$ we have 

$$\sum_{i \in [s]} c_id_i \leq |I|\frac{k^2}{4} + 2|O_1| \cdot e(k-e).$$ 
\end{lemma}

\begin{proof}
We find an upper bound on the maximum of $\sum_{i \in [s]} c_id_i$ under the constraints above. The maximum exists because $\sum_{i \in [s]} c_id_i$ is continuous and the constrained space is compact with respect to $c_i,d_i$. We may assume $c_i + d_i = k$ for all $i \in [s]$, because if $c_j + d_j < k$ for some $j \in [s]$, then we may increase both $c_j,d_j$ so that $\sum_{i \in [s]} c_id_i$ increases. Thus we can just consider maximizing  

$$f(c_1,c_2,\ldots,c_{s}) = \sum_{i \in [s]} c_i(k-c_i)$$

\noindent
and the constraint $\sum_{i \in I \sqcup O_2} d_i \geq 1$ becomes

$$\sum_{i \in I \sqcup O_2} k-c_i  \geq 1.$$

We now show we can assume $\sum_{i \in I \sqcup O_1} c_i  = 1$ and $\sum_{i \in I \sqcup O_2} k-c_i  = 1$. We first note that the derivative of $g(x) = x(k-x)$ is $g'(x) = k-2x$, and so $g(x)$ is maximized at $x = \frac{k}{2}$. Furthermore, $g''(x) = -2$ so that if $x$ approaches $\frac{k}{2}$ we know $g(x)$ increases. If we had $\sum_{i \in I \sqcup O_1} c_i  > 1$, this would imply for some $j \in I \sqcup O_1$ we had $c_j > \frac{1}{|I| + |O_1|} > \frac{k}{2}$, where the last inequality comes from $k < \frac{2}{|I|+|O_1|}$. Thus we may decrease $c_j$ so that $f(c_1,c_2,\ldots,c_{s}) = \sum_{i \in [s]} c_i(k-c_i)$ attains a larger value while still maintaining the constraints. So let $\sum_{i \in I \sqcup O_1} c_i  = 1$ be the new constraint in place of $\sum_{i \in I \sqcup O_1} c_i  \geq 1$. \\

Now we show we may assume $\sum_{i \in I \sqcup O_2} k-c_i = 1$. So assume otherwise that $\sum_{i \in I \sqcup O_2} k-c_i > 1$. Note if for some $j \in O_2$ we had $k -  c_j > \frac{k}{2}$, then we may increase $c_j$ so that $f(c_1,c_2,\ldots,c_{s}) = \sum_{i \in [s]} c_i(k-c_i)$ attains a larger value while still maintaining the constraints. So assume for all $i \in O_2$, we have $k-c_i \leq \frac{k}{2}$. If for some $j \in O_2$ we had $k-c_j < \frac{k}{2}$, then we may decrease $c_j$ so that $f(c_1,c_2,\ldots,c_{s}) = \sum_{i \in [s]} c_i(k-c_i)$ attains a larger value while still maintaining the constraints. So assume $k-c_i = \frac{k}{2}$ for all $i \in O_2$. Thus $\sum_{i \in O_2} k-c_i = |O_2|\frac{k}{2}$, so along with $\sum_{i \in I \sqcup O_2} k-c_i > 1$ we get $\sum_{i \in I} k-c_i > 1 - | O_2|\frac{k}{2}$, and so 

$$\sum_{i \in I} c_i < |I|k + |O_2|\frac{k}{2} - 1 = (2|I| + |O_2|)\frac{k}{2}-1. $$ 

Since $\sum_{i \in I \sqcup O_1} c_i = 1$ and $\sum_{i \in I} c_i < (2|I| + |O_2|)\frac{k}{2}-1$, we must have 

$$\sum_{i \in O_1} c_i = 1-\sum_{i \in I} c_i > 2-(2|I| + |O_2|)\frac{k}{2}$$

\noindent
and thus for some $j \in O_1$ we have $c_j > \frac{2}{|O_1|} - \frac{k(2|I|+|O_1|)}{2|O_1|} > \frac{k}{2}$, where the last inequality can be rearranged into $\frac{k}{2} < \frac{1}{|I|+|O_1|}$. From $\sum_{i \in I} c_i <  (2|I| + |O_2|)\frac{k}{2}-1 $ we also get for some $l \in I$ we have $c_l < (2+\frac{|O_2|}{|I|})\frac{k}{2} - \frac{1}{|I|} < \frac{k}{2}$ where the last inequality can also be rearranged into $\frac{k}{2} < \frac{1}{|I|+|O_1|}$. \\

Note we can increase $c_l$ by $\epsilon > 0$ and decrease $c_j$ by $\epsilon$ so that $f(c_1,c_2,\ldots,c_{s}) = \sum_{i \in [s]} c_i(k-c_i)$ attains a larger value while still maintaining the constraints $\sum_{i \in I \sqcup O_1} c_i = 1$ and $\sum_{i \in I \sqcup O_2 } k-c_i \geq 1$. The constraint $\sum_{i \in I \sqcup O_1} c_i = 1$ is preserved because the $\epsilon$'s cancel out. The constraint $\sum_{i \in I \sqcup O_2 } k-c_i \geq 1$ can be preserved for small enough $\epsilon$ because by assumption $\sum_{i \in I \sqcup O_2 } k-c_i > 1$. Thus under the assumption $\sum_{i \in I \sqcup O_2 } k-c_i > 1$, we can always find a larger value of $f(c_1,c_2,\ldots,c_{s}) = \sum_{i \in [s]} c_i(k-c_i)$ while maintaining the desired constraints, so we may assume $\sum_{i \in I \sqcup O_2} k-c_i = 1$. \\

We claim under the constraints $\sum_{i \in I \sqcup O_1} c_i = 1$ and $\sum_{i \in I \sqcup O_2 } k-c_i = 1$, the function $f(c_1,c_2,\ldots,c_{s}) = \sum_{i \in [s]} c_i(k-c_i)$ has a maximum. From the constraint $\sum_{i \in I \sqcup O_2 } k-c_i = 1$ we get that $\sum_{i \in I \sqcup O_2} c_i = (|I| + |O_2|)k-1$. Note we have

\begin{align*}
f(c_1,c_2,\ldots,c_{s}) &= \sum_{i \in [s]} c_i(k-c_i) = k \left( \sum_{i \in [s]} c_i \right) - \sum_{i \in [s]} c_i^2 \\
&= k \left( \sum_{i \in I \sqcup O_1} c_i + \sum_{i \in O_2} c_i \right) - \sum_{i \in [s]} c_i^2 = k \left( 1+ \sum_{i \in O_2} c_i \right) - \sum_{i \in [s]} c_i^2.
\end{align*}

\noindent
This shows for any $i \in I \sqcup O_1$, we have 

$$f(c_1,c_2, \ldots, c_i, \ldots, c_s) = f(c_1,c_2, \ldots, |c_i|, \ldots, c_s)$$

\noindent
and analogously 

\begin{align*}
f(c_1,c_2,\ldots,c_{s}) &= \sum_{i \in [s]} c_i(k-c_i) = k \left( \sum_{i \in [s]} c_i \right) - \sum_{i \in [s]} c_i^2 \\
&= k \left( \sum_{i \in O_1} c_i + \sum_{i \in I \sqcup O_2} c_i \right) - \sum_{i \in [s]} c_i^2 \\
&= k \left( \left( \sum_{i \in O_1} c_i \right) +  (|I| + |O_2|)k-1  \right) - \sum_{i \in [s]} c_i^2
\end{align*}

\noindent 
which shows for any $i \in I \sqcup O_2$, we have 

$$f(c_1,c_2, \ldots, c_i, \ldots, c_s) = f(c_1,c_2, \ldots, |c_i|, \ldots, c_s).$$ 

\noindent
Thus $f(c_1,c_2,\ldots,c_s) = f(|c_1|,|c_2|,\ldots,|c_s|)$, so we may assume $c_i \geq 0$ for all $i \in [s]$. Note along with the assumption $c_i \geq 0$ for all $i \in [s]$, the constraints $\sum_{i \in I \sqcup O_1} c_i = 1$ and $\sum_{i \in I \sqcup O_2 } k-c_i = 1$ give rise to a compact space, and since $f$ is continuous we know $f$ attains a maximum in this compact space defined by the assumption $c_i \geq 0$ for all $i \in [s]$ with the constraints $\sum_{i \in I \sqcup O_1} c_i = 1$ and $\sum_{i \in I \sqcup O_2 } k-c_i = 1$, and thus $f$ attains a maximum with the constraints $\sum_{i \in I \sqcup O_1} c_i = 1$ and $\sum_{i \in I \sqcup O_2 } k-c_i = 1$ only. \\

We use Lagrange multipliers to maximize $f(c_1,c_2,\ldots,c_{s}) = \sum_{i \in [s]} c_i(k-c_i)$ under the constraints $\sum_{i \in I \sqcup O_1} c_i = 1$ and $\sum_{i \in I \sqcup O_2 } k-c_i = 1$. We have $\frac{\partial f}{\partial c_i} = k - 2c_i$ for all $i \in [s]$. Define $g_1(c_1,c_2,\ldots,c_s) = \sum_{i \in I \sqcup O_1} c_i$ and $g_2(c_1,c_2,\ldots,c_s) = -(\sum_{i \in I \sqcup O_2 } k-c_i)$. So $\frac{\partial g_1}{\partial c_i} = 1$ if $i \in I \cup O_1$ and $0$ otherwise, and $\frac{\partial g_2}{\partial c_i} = 1$ if $i \in I \cup O_2$ and $0$ otherwise. Thus for some $\lambda_1,\lambda_2 \in \mathbb{R}$ we have 
$k-2c_i = \lambda_1+\lambda_2$ for all $i \in I$, and $k-2c_i = \lambda_1$ for all $i \in O_1$, and $k-2c_i = \lambda_2$ for all $i \in O_2$. This gives us for all $i,j \in I$ we have $c_i = c_j$, for all $i,j \in O_1$ we have $c_i = c_j$, and for all $i,j \in O_2$ we have $c_i = c_j$. \\

Without loss of generality, let $c_1 \in O_1, c_2 \in O_2, c_3 \in I$. We get 

$$\sum_{i \in I \sqcup O_1} c_i = \sum_{i \in I} c_i + \sum_{i \in O_1} c_i = \sum_{i \in I} c_3 + \sum_{i \in O_1} c_1 = |I|c_3 + |O_1|c_1 = 1.$$

\noindent
Likewise

$$\sum_{i \in I \sqcup O_2} k-c_i = \sum_{i \in I}k- c_i + \sum_{i \in O_2}k- c_i = |I|(k-c_3) + |O_2|(k-c_2) = 1.$$

\noindent
Thus we get $|I|c_3 + |O_1|c_1 =|I|(k-c_3) + |O_2|(k-c_2)$, and so 

$$2|I|c_3 = |I|k+|O_1|(k-c_2-c_1)$$

\noindent
which gives us $2c_3 = k+ \frac{|O_1|}{|I|}(k-c_2-c_1)$. Recall also that $k-2c_3 = \lambda_1 + \lambda_2 = k-2c_1+k-2c_2 $ and thus $2c_3 = 2c_1+2c_2-k$. Thus we have 

$$2c_1+2c_2-k = 2c_3 = k+\frac{|O_1|}{|I|}(k-c_2-c_1)$$

\noindent
which rearranges into

$$0 = (\frac{|O_1|}{|I|} + 2)(k-c_2-c_1)$$

\noindent
and thus $c_1+c_2 = k$. \\

Recall $2c_3 = 2c_1+2c_2 - k$, so along with $c_1+c_2 = k$ we get $2c_3 = k$ so that $c_3 = \frac{k}{2}$. Recall also that $|I|c_3 + |O_1|c_1 = 1$, so that 

$$c_1 = \frac{1}{|O_1|}(1-|I|\frac{k}{2}) = \frac{2-|I|k}{2|O_1|} = e.$$

\noindent
Note from $c_1+c_2 = k$ we get $c_1(k-c_1) = c_2(k-c_2)$ Finally, we compute

\begin{align*}
f(c_1,c_2,\ldots,c_s) &= |I|c_3(k-c_3) + |O_1|c_1(k-c_1)+|O_2|c_2(k-c_2) \\
&= |I|\frac{k^2}{4} + 2|O_1| \cdot e(k-e) 
\end{align*}

\noindent
as required.
\end{proof}

\end{document}
